
\myaddRA{\subsection{Dead-Path-Aware DPCM for General \texorpdfstring{$r$}{r}}}
\label{sec:opt}

The previous subsections instantiate DPCM for $r{=}1$, $r{=}2$, and general $r$.
%
For $r{=}1$ and $r{=}2$, the dominant updates are already reduced to counter-only or closed-form operators.
%
The remaining bottleneck appears in the general regime, where conflict-state DPCM may still invoke recursive path splitting.
%
In this subsection, we introduce a stronger journal-level refinement for that regime: \emph{dead-path-aware} maintenance.
%
The idea is to use the maintained conflict state to certify that many affected paths cannot possibly encode any surviving size-$s$ clique, allowing us to discard them before recursion.

\stitle{Key Idea: Budgeted Hitting Set.}
When multiple $r$-cliques are removed from a single clique path $\TPath$, Algorithm~\ref{alg:removeRClique} must determine which subpaths survive.
Before invoking the recursive search, we apply a polynomial-time test that detects \emph{dead paths}---paths that cannot produce any valid subpath of size $\ge s$.
If the test succeeds, we skip the exponential splitting entirely and subtract the remaining support in $O\bigl(\binom{|\TPath|}{r}\bigr)$ time.

\begin{definition}[Conflict Family]
Given a set of removed $r$-cliques $C_{\mathrm{rm}} = \{C_0, \dots, C_{m-1}\}$ on path $\TPath$,
the \emph{conflict family} is $\mathcal{F} = \{F_0, \dots, F_{m-1}\}$ where $F_i = C_i \cap V_p(\TPath)$ is the set of \pivot members of $C_i$.
The \emph{budget} is $B = |V_p(\TPath)| - (s - |V_h(\TPath)|)$: the maximum number of pivots that may be removed.
\end{definition}

Any valid subpath must exclude at least one pivot from each $F_i$ while retaining at least $s - |V_h|$ pivots.
This is precisely a \emph{budgeted hitting set} problem.
Computing the exact minimum hitting set is NP-hard, but the following polynomial-time conditions suffice for most practical instances.

\begin{lemma}[Forced Vertex]\label{lem:forced_vertex}
If a vertex $v$ belongs to every $r$-clique on $\TPath$ (i.e., its conflict degree equals $\binom{|\TPath|-1}{r-1}$), then $v$ must be removed in every feasible solution.
If $v$ is a \hold, the path is dead.
\end{lemma}

\begin{lemma}[Singleton Propagation]\label{lem:singleton}
If $|F_i| = 1$, say $F_i = \{p\}$, then $p$ must be removed.
Removing $p$ may create new singletons; we propagate iteratively until no singleton remains.
If at any point $|V_p| + |V_h| < s$, the path is dead.
\end{lemma}

\begin{lemma}[Disjoint Packing Lower Bound]\label{lem:packing}
If there exist $t$ pairwise disjoint sets in $\mathcal{F}$, then any feasible solution must remove $\ge t$ pivots.
If $t > B$, the path is dead.
\end{lemma}

\begin{lemma}[Component Decomposition]\label{lem:component_lb}
Let $G_c$ be the graph on $V_p$ where two pivots are adjacent iff they co-occur in some $F_i$.
If $G_c$ has connected components $C_1, \dots, C_k$, then the minimum hitting set decomposes as $\tau(\mathcal{F}) = \sum_j \tau(\mathcal{F}[C_j])$.
Computing the packing lower bound per component and summing yields a bound at least as tight as the global one.
\end{lemma}

These conditions are applied as a preprocessing step before Algorithm~\ref{alg:removeRClique}.
Algorithm~\ref{alg:pathsplit_pruning} summarizes the procedure.


\begin{algorithm}[t]
\small
\caption{\textsc{PrunedUpdateIndex}}
\label{alg:pathsplit_pruning}
\KwIn{clique path $\TPath{} = \{V_H, V_P\}$, $r$-cliques to remove $C_{\mathrm{rm}}$, clique size $s$}
\KwOut{All subpaths after deletion, or $\emptyset$ if dead}
\SetKwFunction{Split}{PathSplit}
\SetKwProg{Fn}{Function}{}{}

$F_i \gets C_i \cap V_P$ for each $C_i \in C_{\mathrm{rm}}$\tcp*{conflict sets}
$B \gets |V_P| - (s - |V_H|)$\tcp*{removal budget}

\tcp{Forced vertex (Lemma~\ref{lem:forced_vertex})}
\ForEach{$v$ with conflict degree $\ge \binom{|\TPath|-1}{r-1}$}{
  \lIf{$v \in V_H$}{\Return $\emptyset$}
  $V_P \gets V_P \setminus \{v\}$; resolve all $F_i \ni v$\;
}

\tcp{Singleton propagation (Lemma~\ref{lem:singleton})}
\Repeat{no change}{
  \ForEach{active $F_i$ with $|F_i| = 1$, say $\{p\}$}{
    $V_P \gets V_P \setminus \{p\}$; resolve all $F_j \ni p$\;
    \lIf{$|V_H| + |V_P| < s$}{\Return $\emptyset$}
  }
}

\tcp{Component packing LB (Lemmas~\ref{lem:packing},~\ref{lem:component_lb})}
Build conflict graph $G_c$; find components $C_1, \dots, C_k$\;
$\mathit{LB} \gets \sum_{j=1}^{k} \textsc{GreedyPacking}(\{F_i \mid F_i \subseteq C_j\})$\;
\lIf{$\mathit{LB} > B$}{\Return $\emptyset$}

\Return \Split{$V_H, V_P$} with reduced $\mathcal{F}$\tcp*{Alg.~\ref{alg:removeRClique}}
\end{algorithm}



\stitle{Effectiveness.}
Our experiments show that the above tests prune the vast majority of affected paths on representative datasets.
At $s = 20$, up to 97\% of paths are detected as dead, reducing the number of paths that enter the exponential splitting (Algorithm~\ref{alg:removeRClique}) by up to 47$\times$.
This is the primary reason our method completes on parameter ranges where the baseline times out.
